\documentclass[
  12pt,			% Schriftgröße
	a4paper,			% Papierformat
	oneside,			% einseitiges Dokument
	titlepage,			% es wird eine Titelseite verwendet
	parskip=full,		% Abstand zwischen Absätzen (volle Zeile)=full halbe zeille = half
	headings=normal,	% Größe der Überschriften verkleinern
	BCOR=10mm,		% Bindekorrektur
	final				% Status des Dokuments (final/draft)
]{scrreprt}

\usepackage{graphicx}
%\usepackage{wasysym}
\usepackage{pstricks,pst-node,pst-tree} %Für Bäume

%%Smiley ;-)
\newcommand{\smiley}{
\begin{pspicture}(0.5,0.35)
	\pscircle(0.25,0.15){0.2}
	\qdisk(0.19,0.20){0.03}
	\qdisk(0.31,0.20){0.03}
	\psarc(0.25,0.15){0.1}{210}{330}
\end{pspicture}
}

\usepackage[latin1]{inputenc}
\usepackage[T1]{fontenc}
\usepackage{setspace}
\usepackage[
	automark,		% Kapitelangaben in Kopfzeile automatisch erstellen
	headsepline,		% Trennlinie unter Kopfzeile
	ilines			% Trennlinie linksbündig ausrichten
]{scrpage2}
\usepackage{lmodern}	%Schriftbild in Adobe PDf sonst unleserlich
 
%Formatierung geht ab hier los:
 
\setkomafont{sectioning}{\normalcolor\bfseries} % Ändert Überschriftschriftart (nun mit Serifen)
\onehalfspacing																	%Zeilenabstand 1,5 fach
\KOMAoptions{DIV=last}		%Satzspiegel neu berechnen wegen 1,5 fach Zeilenabstand
 
% Kopf- und Fußzeilen ------------------------------------------------------
\pagestyle{scrheadings}
 
% Kopf- und Fußzeile auch auf Kapitelanfangsseiten -------------------------
\renewcommand*{\chapterpagestyle}{scrheadings}
 
% Schriftform der Kopfzeile ------------------------------------------------
\renewcommand{\headfont}{\normalfont}
 
% Kopfzeile ----------------------------------------------------------------
\ihead{\textit{\headmark}}
\chead{}
\chead{}
\setlength{\headheight}{21mm} % Höhe der Kopfzeile
\setheadwidth[0pt]{textwithmarginpar} % Kopfzeile über den Text hinaus verbreitern
\setheadsepline[text]{0.4pt} % Trennlinie unter Kopfzeile
 
% Fußzeile -----------------------------------------------------------------
%\ifoot{\copyright\ \autor}
\cfoot{}
\ofoot{\pagemark}

\title{Automatentheorie und formale Sprachen}
\subject{Script zur Lehrveranstaltung}
\author{Florian Schleich}
\subtitle{Prof. Annika Wagner}
\publishers{\includegraphics[scale=0.5]{images/HS-Fulda-logo.eps}}

\begin{document}
\maketitle

\chapter{Grammatiken}
\pstree[levelsep=8ex]{\Tcircle{R} }{
\Tcircle{D1}
\pstree{ \Tcircle{D2} }{
\Tcircle{D2.1}
\Tcircle{D2.2}
}
\Tcircle{D3}
\Tcircle{D4}
}

\section{Symbole}
Symbole sind alle verwendbaren Zeichen eines Alphabets.

\section{Alphabet}
Ein Alphabet, abgek�rzt $\Sigma$, ist eine endliche Menge von einzelnen Symbolen. Ein Alphabet wird als endliche Menge von Symbolen definiert.

$\Sigma = \{\smiley, a, Anton\}$

Neben einzelnen Buchstaben k�nnen auch W�rter im Sinne der deutschen Sprache verwendet werden. Es k�nnen aber auch andere Symbole (\smiley) verwendet werden. Bei Verwendung von W�rtern als Symbole eines Alphabets ist aber zu beachten, dass bei der sp�teren Verwendung die einzelnen Buchstaben oder Teile des Wortes nicht als Symbole genommen werden d�rfen, sondern das gesamte Wort als ein Symbol bezeichnet wird. Im hier verwendeten Fall darf also das gro�e $A$ nicht als Symbol verwendet werden darf. ebenso d�rfen auch keine Wortteile wie $ton$ nicht verwendet werden. Ein zul�ssiges Symbol ist allerdings das gesamte Wort $Anton$.

\section{W�rter}
W�rter werden aus Kombination von einzelnen Symbolen gebildet. Die bildbaren W�rter werden in einer Menge $\Sigma*$ gesammelt. Wird die Menge aller W�rter gebildet muss jetzt explizit das leere Wort $\epsilon$ angegeben werden. Auf Grund der in der Grammatik verwendeten Regeln wird die Menge der m�glichen W�rter eingeschr�nkt. In unserem Beispiel kann die Menge aller W�rter so aussehen:

$\Sigma* = \{\epsilon, \smiley, a, Anton, \smiley a, a Anton, \smiley Anton, a \smiley, Anton a, Anton \smiley, \smiley \smiley, ...\}$

Die Menge wird gebildet, in dem erst das leere Wort, dann alle einstelligen W�rter, alle zweistelligen W�rter, alle dreistelligen W�rter usw. in die Menge integriert werden. Die Menge der W�rter ist unendlich, da die L�nge der W�rter nicht bestimmt ist.

Die L�nge eines Wortes $|w|$ bezeichnet die Anzahl der f�r das Wort verwendeten Symbole.

\begin{figure}[h]
	$|\epsilon| = 0$
	
	$|a| = 1$
	
	$|Anton| = 1$
	
	$|aa| = 2$
	\caption{Beispiele f�r Wortl�ngen}
\end{figure}

\section{Grammatiken}
Eine Grammatik $G$ wird definiert durch:
\begin{itemize}
	\item eine Menge von Variablen $V=\{S,U\}$, die bei der Ableitung verschwinden,
	\item eine Menge von Terminalsymbolen $\Sigma$, gelichbedeutend mit dem definierten Alphabet
	\item ein Startsymbol $S \in V$
	\item und eine Menge von Regeln $P$
\end{itemize}

Formal wird eine Grammatik also beschrieben als:

$G = (V, \Sigma, S, P)$

\subsection{Regeln}
Eine Regel besteht aus einer linken und einer rechten Regelseite. Bei der Ableitung von Regeln wird der Ausdruck auf der linken Seite durch den Ausdruck auf der rechten Seite ersetzt.

Bei der Regel $S \rightarrow Anton$ wird die Variable $S$ durch das Symbol $Anton$ ersetzt. 

\subsection{Ableitung}
Bei der Ableitung von Regeln beginnt man beim Startsymbol und ersetzt das Startsymbol durch den entsprechenden Ausdruck einer Regel, die auf das Startsymbol passt. Zwischen zwei Ableitungsschritten kommt dann ein Doppelpfeil ($\Rightarrow$).

\begin{figure}[h]
$S \rightarrow Anton$ \textbf{(Regel)}

$S \Rightarrow Anton$ \textbf{(Ableitung)}
\caption{Einfaches Beispiel f�r eine Ableitung}
\end{figure}

\subsection{Regelalternativen}
Regeln k�nnen auch Alternativen enthalten, um mehr M�glichkeiten zu bieten W�rter zu erzeugen.
Alternativen werden auf der rechten Regelseite durch ein Pipe-Symbol ($|$) eingeleitet.
Eine Regel mit Alternativen k�nnte folgenderma�en aussehen:

$S \rightarrow Anton | Anton \smiley | Anton a$

Dadurch w�ren die W�rter $Anton$, $Anton \smiley$ und $Anton a$ f�r die Variable $S$ m�glich.
\chapter{Endliche Automaten}
Endliche Automaten werden dazu genutzt um zu pr�fen ob ein Wort auf eine bestimmte Grammatik passt. Diese Eigenschaft nutzt ein Parser um den Quellcode auf seine syntaktische Richtigkeit zu pr�fen. Um aus einer Grammatik einen Automat herzuleiten, ist es wichtig zu beachten, dass die Grammatik regul�r ist.
\section{Regul�re Grammatiken und deren Verabreitung}
Endliche Automaten werden als Zustandsdiagramm dargestellt. Ein Zustandsdiagramm f�r eine Lampe l�sst sich in Abbildung 
\ref{fig:zustandLampe} finden.

\begin{figure}[h]
	\psmatrix[mnode=circle,colsep=1]
	 an & & aus & \psdot(0,0)
	\endpsmatrix
	\psset{shortput=nab,arrows=->,labelsep=3pt}
	\small
	\ncarc[arcangle=12]{1,1}{1,3}^{\scriptsize{dr�cken}}
	\ncarc[arcangle=12]{1,3}{1,1}^{\scriptsize{dr�cken}}
	\ncarc[arcangle=12]{1,4}{1,3}

	\caption{Zustandsdiagramm f�r eine Lampe}
	\label{fig:zustandLampe}
\end{figure}

Mit jedem Druck auf eine Taste wird der Status der Lampe ge�ndert. Der Startzustand ist dabei "aus".

Ein Automat f�r die schon vorher besprochene Grammatik: "`L�nge des Wortes gerade"', k�nnte nach einer kurzen �berlegung wie in Abbildung \ref{fig:zustandIntuGerade} aussehen.

Die zugeh�rige formale Grammatik findet sich in Grafik \ref{fig:grammatikGerade}

\begin{figure}[h]
		$\Sigma = \{0,1\}$
		
		$S \rightarrow 0T | 1T | \epsilon$
		
		$T \rightarrow 0S | 1S$	

	\label{fig:grammatikGerade}
	\caption{Grammatik f�r W�rter mit gerader L�nge}
\end{figure}

\begin{figure}[h]
	\psmatrix[mnode=circle,colsep=1]
	 . \\
	 gerade & & ungerade \\
	 \psdot(0,0)
	\endpsmatrix
	\psset{shortput=nab,arrows=->,labelsep=3pt}
	\small
	\ncarc[arcangle=12]{1,1}{2,1}
	\ncarc[arcangle=12]{2,1}{2,3}^{\scriptsize{0,1}}
	\ncarc[arcangle=12]{2,3}{2,1}^{\scriptsize{0,1}}
	\ncarc[arcangle=12]{2,1}{3,1}	

	\caption{Intuitiver Automat f�r gerade W�rter}
	\label{fig:zustandIntuGerade}
\end{figure}

Dieser Automat ist leicht zu verstehen. Wir kommen durch den Startzustand in den Zustand "gerade". Mit einem Eingabesymbol (1 oder 0) kommen wir in den Zustand "ungerade", da das Wort jetzt die L�nge 1 hat und damit ungerade ist.F�gen wir ein weiteres Eingabesymbol an, kommen wir wieder in den Zustand "gerade". Da dieser Zustand auch gleichzeitig Endzustand ist, wird das Wort mit der L�nge 2 erkannt.

\section{Erzeugen von endlichen Automaten}
Nun stellt sich aber die Frage, wie man aus einer regul�ren Grammatik einen endlichen Automaten ableitet. F�r diese Problemstellung, gibt es ein Kochrezept, welches nun kurz vorgestellt wird.

\subsection{Schritt 1: Zust�nde definieren}
Zust�nde in Automaten sind gleichbedeutend mit den verwendeten Variablen der Grammatik. F�r jede Variable, die in der Grammatik definiert wurde, gibt es auch ein speziellen Zustand. In unserem Fall wurden die Variablen $S$  und $T$ verwendet.

\begin{figure}[h]
	\psmatrix[mnode=circle,colsep=1]
	 S & & T
	\endpsmatrix

	\caption{Automat nach 1. Schritt}
\end{figure}

\subsection{Schritt 2: Startzustand einzeichnen}
Der Startzustand ergibt sich aus der formalen Definition der Grammatik. Da Zust�nde gleichbedeutend mit Variablen sind, ist auch der Startzustand gleichbedeutend mit der Startvariable.

Um einen Startzustand einzuzeichen, wird ein dickerer Punkt mittels Pfeil ohne Beschriftung auf den Startzustand gezeichnet.

\begin{figure}[h]
	\psmatrix[mnode=circle,colsep=1]
	 . \\
	 S & & T
	\endpsmatrix
	\psset{shortput=nab,arrows=->,labelsep=3pt}
	\small
	\ncarc[arcangle=12]{1,1}{2,1}

	\caption{Automat nach Schritt 2}
\end{figure}

\subsection{Schritt 3: Endzust�nde einzeichnen}
Endzust�nde sind all diejenigen, bei denen die Verarbeitung eines Wortes aufh�ren kann. Klassischerweise enth�lt ein Endzustand das leere Wort ($\epsilon$). Allerdings kann die Verarbeitung auch bei einem einzelnen Terminalsymbol beendet werden. Dazu aber an sp�terer Stelle mehr. Ein Endzustand wird durch einen Punkt mit umschlossenen Kreis dargestellt. Auf diesen Kreis zeigt dann ausgehend von einem Endzustand ein Pfeil. Es kann durchaus vorkommen, dass es mehrere Endzust�nge gibt, wohingegen nur ein Startzustand existieren darf.

\begin{figure}[h]
	\psmatrix[mnode=circle,colsep=1]
	 . \\
	 S & & T \\
	 \psdot(0,0)
	\endpsmatrix
	\psset{shortput=nab,arrows=->,labelsep=3pt}
	\small
	\ncarc[arcangle=12]{1,1}{2,1}
	\ncarc[arcangle=12]{2,1}{3,1}	

	\caption{Automat nach dem 3. Schritt}
\end{figure}

\subsection{Schritt 4: Kanten einzeichnen}
Eine Kante oder Transition in einem Automat, wird durch eine Regel erzeugt. Da die Regeln in einer regul�ren Grammatik immer einem einfachen Shema folgen, ist die Erstellung von Kanten kein Problem. Kanten werden durch Pfeile auf andere Zust�nde dargestellt. Der linke Regelteil definiert den Anfang des Pfeils. Im rechten Regelteil stehen immer ein Termianlsymbol gefolgt von einer Variable. Die Variable im rechten Regelteil gibt an, auf welchen Zustand die Kante zeigen soll. Das Terminalsymbol der rechten Regelseite wird an die Kante geschrieben.

\begin{figure}[h]
	\psmatrix[mnode=circle,colsep=1]
	 . \\
	 S & & T \\
	 \psdot(0,0)
	\endpsmatrix
	\psset{shortput=nab,arrows=->,labelsep=3pt}
	\small
	\ncarc[arcangle=12]{1,1}{2,1}
	\ncarc[arcangle=12]{2,1}{2,3}^{\scriptsize{0,1}}
	\ncarc[arcangle=12]{2,3}{2,1}^{\scriptsize{0,1}}
	\ncarc[arcangle=12]{2,1}{3,1}	

	\caption{Automat nach Schritt 4}
\end{figure}

\subsection{Sonderfall bei regul�ren Grammatiken}
Reg�l�re Grammatiken lassen auch Regeln nach dem Schema $s \rightarrow a$ zu. Diese Regeln k�nnen nach dem bisher besprochenen Schema nicht in eine Kante f�r einen Automaten umgewandelt werden. Um diese Art von Regeln zu behandeln, ist ein kleiner Trick von n�ten. Man erzeugt einfach einen neuen Zustand. Dieser Zustand wird durch die Variable auf der linken Regelseite und das Terminalsymbol benannt ($S_a$). Die Kante die zu diesem Zustand f�hrt wird mit dem Terminalsymbol beschriftet. Der neue Zustand wird dann selbst zu einen Endzustand. Im folgender Grafik nehmen wir an, dass eine Regel $S \rightarrow a$ existiert.

\begin{figure}[h]
	\psmatrix[mnode=circle,colsep=1]
	 . & $S_a$ &\psdot(0,0)\\
	 S & & T \\
	 \psdot(0,0)
	\endpsmatrix
	\psset{shortput=nab,arrows=->,labelsep=3pt}
	\small
	\ncarc[arcangle=12]{1,1}{2,1}
	\ncarc[arcangle=12]{2,1}{2,3}^{\scriptsize{0,1}}
	\ncarc[arcangle=12]{2,3}{2,1}^{\scriptsize{0,1}}
	\ncarc[arcangle=12]{2,1}{3,1}
	\ncarc[arcangle=12]{2,1}{1,2}^{\scriptsize{a}}
	\ncarc[arcangle=12]{1,2}{1,3}

	\caption{Automat mit zus�tzlichem Zustand $S_a$}
\end{figure}

\section{Determinismus}
In der Automaten Therorie wird zwischen deterministischen und nicht deterministischen Automaten unterschieden. Ein nicht deterministischer Automat hat mindestens eine der folgenden Eigenschaften.

\paragraph{Uneindeutigkeit}
Ein Automat ist uneindeutig, wenn es f�r einen Zustand in einem Terminalsymbol mehrere Kanten und damit M�glichkeiten weiterzugehen gibt.

\paragraph{Unterspezifiziert}
Umgekehrt kann ein Automat unterspezifiziert sein, wenn es f�r einen Zustand und ein Terminalsymbol keine ausgehende Kante gibt.

Ist auch nur eine dieser Eigenschaften erf�llt, wird der Automat als nicht deterministisch eingestuft.

\section{Suchb�ume}
Nicht deterministische Automaten sind nicht empfehlenswert. Um ein Wort zu parsen, was die Hauptaufgabe von Automaten ist, m�ssen bei nicht deterministischen Suchb�ume eingesetzt werden. Um bei einem Suchbaum zu einem Ergebnis zu kommen, m�ssen m�glicherweise alle �ste durchsucht werden. Dies begr�ndet sich dadurch, dass bei nicht deterministischen Automaten mehrere M�glichkeiten existieren, ein Wort zu erkennen.

\begin{figure}[h]
		$\Sigma = \{a,b\}$
		
		$S \rightarrow aS | aT | a | b$
		
		$T \rightarrow bT | \epsilon$	

	\label{fig:grammatikND}
	\caption{Nicht deterministische Grammatik}
\end{figure}

\begin{figure}[h]
	\psmatrix[mnode=circle,colsep=0.5,rowsep=0.5]
	 .\\
	 S & & T \\
	 $S_a$ & $S_b$ & \psdot(0,0)\\ 
	 \psdot(0,0) & \psdot(0,0)
	\endpsmatrix
	\psset{shortput=nab,arrows=->,labelsep=3pt}
	\small
	\ncarc[arcangle=12]{1,1}{2,1}
	\nccurve[angleA=100,angleB=180,offsetB=0,ncurvA=6,ncurvB=6]{2,1}{2,1}^{\scriptsize{a}}
	\ncarc[arcangle=12]{2,1}{2,3}^{\scriptsize{a}}
	\ncarc[arcangle=12]{2,1}{3,1}^{\scriptsize{a}}
	\ncarc[arcangle=12]{2,1}{3,2}^{\scriptsize{b}}
	\nccurve[angleA=80,angleB=0,offsetB=0,ncurvA=6,ncurvB=6]{2,3}{2,3}^{\scriptsize{b}}
	\ncarc[arcangle=12]{2,3}{3,3}
	\ncarc[arcangle=12]{3,1}{4,1}
	\ncarc[arcangle=12]{3,2}{4,2}
	

	\caption{Nicht deterministischer Automat}
\end{figure}

\begin{figure}[h]
\pstree[levelsep=8ex]{\Tcircle{S}}{
	\pstree{ \Tcircle{$S_a$} }{
			\Tcircle{\blitz}
	}
	\pstree{ \Tcircle{T} }{
			\Tcircle{\blitz}
	}
	\pstree{ \Tcircle{S} }{
		\pstree{ \Tcircle{$S_a$} }{
				\Tcircle{\blitz}
		}
		\pstree{ \Tcircle{T} }{
				\Tcircle{T}
		}
		\pstree{ \Tcircle{S} }{
				\Tcircle{$S_b$}
		}
	}
}
\rput(1, -0.5){\psframebox[linewidth=0]{a}}
\rput(1, -2){\psframebox[linewidth=0]{a}}
\rput(1, -3.5){\psframebox[linewidth=0]{b}}
\caption{Suchbaum f�r das Wort "aab"}
\end{figure}

Nicht deterministische Automaten m�ssen wie eingangs schon gesagt �ber Suchb�ume verarbeitet werden. Dadurch, dass es mehrere M�glichkeiten gibt zu einem Ergebnis zu gelangen, muss der Suchbaum durch das Backtracking Verfahren traversiert werden. Der Aufwand beim Backtracking kann quadratisch zum zu parsenden Wort sein. Je nachdem wie schnell ein Ast gefunden wird, der das Wort vollst�ndig erkennt. Bei deterministischen Automaten ist der Aufwand hingegen immer linear. Es werden genau so viele Schritte ben�tigt um das Wort zu erkennen, wie das Wort lang ist. Deterministische Automaten parsen wesentlich schneller, als nicht deterministische.

\section{Lazy Evaluation Verfahren}
Wir haben im letzten Abschnitt festgestellt, dass nicht deterministische Automaten nicht optimal sind. Ziel ist es einen dterministischen Automaten herzuleiten. Im folgenden wird ein Verfahren aufgezeigt, dass es erm�glicht aus einem nicht deterministischen einen deterministischen Automaten herzustellen.

\begin{figure}[h]
	\psmatrix[mnode=circle,colsep=2,rowsep=0.5]
	 . & & &  \psdot(0,0) & $B_0$ & \psdot(0,0)\\
	 & S & A & B & $B_1$ & \psdot(0,0)\\
	 & & & & $B_\#$ & \psdot(0,0)
	\endpsmatrix
	\psset{shortput=nab,arrows=->,labelsep=3pt}
	\small
	\nccurve[angleA=190,angleB=270,offsetB=0,ncurvA=6,ncurvB=6]{2,2}{2,2}^{\scriptsize{0,1}}
	\nccurve[angleA=190,angleB=270,offsetB=0,ncurvA=7,ncurvB=7]{2,4}{2,4}^{\scriptsize{0,1,\#}}
	\ncarc[arcangle=12]{1,1}{2,2}
	\ncarc[arcangle=12]{2,2}{2,3}^{\scriptsize{0,1}}
	\ncarc[arcangle=12]{2,3}{2,4}^{\scriptsize{\#}}
	\ncarc[arcangle=-24]{2,2}{2,4}^{\scriptsize{\#}}
	\ncarc[arcangle=12]{2,4}{1,4}
	\ncarc[arcangle=12]{2,4}{1,5}^{\scriptsize{0}}
	\ncarc[arcangle=12]{2,4}{2,5}^{\scriptsize{1}}
	\ncarc[arcangle=12]{2,4}{3,5}^{\scriptsize{\#}}
	\ncarc[arcangle=12]{1,5}{1,6}
	\ncarc[arcangle=12]{2,5}{2,6}
	\ncarc[arcangle=12]{3,5}{3,6}

	\caption{Beispiel f�r einen nicht deterministischen Automaten}
\end{figure}

\begin{figure}[h]
\pstree[levelsep=8ex]{\Tcircle{S}}{
	\pstree{ \Tcircle{S} }{
			\pstree{ \Tcircle{S} }{
					\pstree{\Tcircle{B}} {
						\pstree{\Tcircle{B}}{}
						\pstree{\Tcircle{$B_0$}}{}
					}
			}
			\pstree{ \Tcircle{A} }{
					\pstree{\Tcircle{B}} {
						\pstree{\Tcircle{B}}{}
						\pstree{\Tcircle{$B_0$}}{}
					}
			}
	}
	\pstree{ \Tcircle{A} }{
			\Tcircle{\blitz}
	}
}
\rput(1, -0.5){\psframebox[linewidth=0]{0}}
\rput(1, -2){\psframebox[linewidth=0]{0}}
\rput(1, -3.5){\psframebox[linewidth=0]{\#}}
\rput(1, -5){\psframebox[linewidth=0]{0}}
\caption{Zugeh�riger Suchbaum f�r das Wort "00\#0"}
\end{figure}

Wir werden den nachweislich nicht deterministischen Automaten mit Hilfe des Lazy Evaluation Verfahren einen equivalenten deterministischen Automaten bilden.

\subsection{Schritt 1: Startzustand}
Der Startzustand des NEA (nicht deterministischer endlicher Automat) wird zum neuen Startzustand des DEA (deterministischer endlicher Automat).

\begin{figure}[h]
	\psmatrix[mnode=circle,colsep=2,rowsep=0.5]
	 . \\
	 & $\{S\}$
	\endpsmatrix
	\psset{shortput=nab,arrows=->,labelsep=3pt}
	\small
	\ncarc[arcangle=12]{1,1}{2,2}

	\caption{Automat nach dem ersten Schritt}
\end{figure}

\subsection{Schritt 2: Folgezustand f�r Startzustand}
Wir nehmen nun den Startzustand und �berpr�fen f�r jedes Eingabesymbol (Terminalsymbol), welche Folgezust�nde es im NEA gibt. Diese Folgezust�nde werden als Menge zusammengefasst und als einen eigenen Zustand im DEA erzeugt.

F�r den Startzustand ergeben sich folgende Folgezust�nde:

\begin{figure}[h]
	$S \stackrel{0}{\longrightarrow} \{S, A\}$
	
	$S \stackrel{1}{\longrightarrow} \{S, A\}$
	
	$S \stackrel{\#}{\longrightarrow} \{B\}$
\end{figure}

Da f�r die Eingabesymbole 0 und 1 die Folgezust�nde gleich sind, werden hierf�r keine neuen Zust�nde eingezeichnet. Es wrd lediglich eine neue Kante mit dem Eingabesymbol eingezeichnet.

\begin{figure}[h]
	\psmatrix[mnode=circle,colsep=2,rowsep=0.5]
	 . & $\{S,A\}$\\
	 & $\{S\}$ & $\{B\}$
	\endpsmatrix
	\psset{shortput=nab,arrows=->,labelsep=3pt}
	\small
	\ncarc[arcangle=12]{1,1}{2,2}
	\ncarc[arcangle=12]{2,2}{1,2}^{\scriptsize{0,1}}
	\ncarc[arcangle=12]{2,2}{2,3}^{\scriptsize{\#}}

	\caption{Automat nach dem zweiten Schritt}
\end{figure}

\subsection{Schritt 3: Folgezustand f�r weitere Zust�nde}
Nun nehmen wir die eben eingezeichneten Folgezust�nde des Startzustandes und errechnen wieder die Folgezust�nde f�r alle Eingabesymbole. Errechnet wird ein Folgezustand, indem wir die Folgezust�nde f�r das entsprechende Eingabesymbol f�r alle Elemente der Menge im Zustand in eine neue Menge vereinigen.

Die Folgezust�nde f�r den Zustand $\{S,A\}$ werden berechnet:

\begin{figure}[h]
	$S \stackrel{0}{\longrightarrow} \{S, A\}$
	
	$A \stackrel{0}{\longrightarrow} \{\}$
	
	(vereinigen: ) $\{S,A\} \cup \{\} = \{S,A\}$
\end{figure}

Der Folgezustand von $\{S,A\}$ mit dem Eingabesymbol 0 ist also $\{S,A\}$. Wir zeichenen also einen Verweis auf sich selbst. Alnalog zu diesem Schema werden auch die Folgezust�nde f�r die anderen Eingabesymbole (1 und \#) berechnet und eingezeichnet.

Dieser Schritt wird solange wiederholt, bis der Automat deterministisch ist. Deterministisch ist ein Automat, wenn er f�r alle Eingabesymbole genau eine ausgehende Kante hat.

\begin{figure}[h]
	\psmatrix[mnode=circle,colsep=1,rowsep=0.5]
	 .\\
	 & $\{S\}$ & $\{S,A\}$\\
	 & & $\{B\}$\\
	 & $\{B,B_0\}$ & $\{B,B_1\}$ & $\{B,B_\#\}$\\
	 \\
	\endpsmatrix
	\psset{shortput=nab,arrows=->,labelsep=3pt}
	\small
	\ncarc[arcangle=12]{1,1}{2,2}
	\ncarc[arcangle=12]{2,2}{2,3}^{\scriptsize{0,1}}
	\ncarc[arcangle=12]{2,2}{3,3}^{\scriptsize{\#}}
	\nccurve[angleA=90,angleB=0,offsetB=0,ncurvA=6,ncurvB=6]{2,3}{2,3}^{\scriptsize{0,1}}
	\ncarc[arcangle=12]{2,3}{3,3}^{\scriptsize{\#}}
	\ncarc[arcangle=12]{3,3}{4,2}^{\scriptsize{0}}
	\ncarc[arcangle=12]{3,3}{4,3}^{\scriptsize{1}}
	\ncarc[arcangle=12]{3,3}{4,4}^{\scriptsize{\#}}
	\ncarc[arcangle=12]{4,2}{4,3}^{\scriptsize{1}}
	\ncarc[arcangle=12]{4,3}{4,4}^{\scriptsize{\#}}
	\ncarc[arcangle=12]{4,4}{4,3}^{\scriptsize{1}}
	\ncarc[arcangle=12]{4,3}{4,2}^{\scriptsize{0}}
	\nccurve[angleA=250,angleB=250,offsetB=0,ncurvA=0.25,ncurvB=0.25]{4,2}{4,4}^{\scriptsize{\#}}
	\nccurve[angleA=280,angleB=280,offsetB=0,ncurvA=0.5,ncurvB=0.5]{4,4}{4,2}^{\scriptsize{0}}

	\caption{Automat nach dem dritten Schritt}
\end{figure}

\subsection{Schritt 4: Endzust�nde}
Jeder Zustand des neu angelegten DEA, der mindestens ein Element enth�lt, dass im NEA auch schon ein Endzustand war, wird zu einem neuen Endzustand.

\begin{figure}[h]
	\psmatrix[mnode=circle,colsep=1,rowsep=2]
	 .\\
	 & $\{S\}$ & $\{S,A\}$\\
	 & & $\{B\}$ & \psdot(0,0)\\
	 & $\{B,B_0\}$ & $\{B,B_1\}$ & $\{B,B_\#\}$\\
	 & \psdot(0,0) & \psdot(0,0) & \psdot(0,0) \\
	\endpsmatrix
	\psset{shortput=nab,arrows=->,labelsep=3pt}
	\small
	\ncarc[arcangle=12]{1,1}{2,2}
	\ncarc[arcangle=12]{2,2}{2,3}^{\scriptsize{0,1}}
	\ncarc[arcangle=12]{2,2}{3,3}^{\scriptsize{\#}}
	\nccurve[angleA=90,angleB=0,offsetB=0,ncurvA=6,ncurvB=6]{2,3}{2,3}^{\scriptsize{0,1}}
	\ncarc[arcangle=12]{2,3}{3,3}^{\scriptsize{\#}}
	\ncarc[arcangle=12]{3,3}{4,2}^{\scriptsize{0}}
	\ncarc[arcangle=12]{3,3}{4,3}^{\scriptsize{1}}
	\ncarc[arcangle=12]{3,3}{4,4}^{\scriptsize{\#}}
	\ncarc[arcangle=12]{4,2}{4,3}^{\scriptsize{1}}
	\ncarc[arcangle=12]{4,3}{4,4}^{\scriptsize{\#}}
	\ncarc[arcangle=12]{4,4}{4,3}^{\scriptsize{1}}
	\ncarc[arcangle=12]{4,3}{4,2}^{\scriptsize{0}}
	\nccurve[angleA=250,angleB=250,offsetB=0,ncurvA=0.25,ncurvB=0.25]{4,2}{4,4}^{\scriptsize{\#}}
	\nccurve[angleA=280,angleB=280,offsetB=0,ncurvA=0.5,ncurvB=0.5]{4,4}{4,2}^{\scriptsize{0}}
	\ncarc[arcangle=12]{3,3}{3,4}
	\ncarc[arcangle=-12]{4,2}{5,2}
	\ncarc[arcangle=0]{4,3}{5,3}
	\ncarc[arcangle=12]{4,4}{5,4}

	\caption{Automat nach dem dritten Schritt}
\end{figure}

Nach der Optimierung nach dem Lazy Evaluation Verfahren wird f�r das Wort 00\#0 folgender Suchbaum erzeugt.

\begin{figure}[h]
	\psmatrix[mnode=circle,colsep=1,rowsep=2]
	 $\{S\}$ & $\{S, A\}$ & $\{S, A\}$ & $\{B\}$ & $\{B, B_0\}$
	\endpsmatrix
	\psset{shortput=nab,arrows=->,labelsep=3pt}
	\small
	
	\ncarc[arcangle=0]{1,1}{1,2}^{\scriptsize{0}}
	\ncarc[arcangle=0]{1,2}{1,3}^{\scriptsize{0}}
	\ncarc[arcangle=0]{1,3}{1,4}^{\scriptsize{\#}}
	\ncarc[arcangle=0]{1,4}{1,5}^{\scriptsize{0}}
\end{figure}

Wir erkennen, dass der zugeh�rige Suchbaum um ein vielfaches geschrumpft ist und zu dem auch linearisiert wurde. Das  bringt enorme Geschwindigkeitsvorteile.
\end{document}